
\section{Introduction}

VALOR is a modern framework for neutrino oscillation analysis, originally established within the T2K experiment in 2010. Its name—an acronym of the founding institutions (VALencia, Oxford, and Rutherford)—has been retained, even as the group has grown significantly in scope and membership.

Since its inception, the VALOR group has carried out approximately 20 major iterations of T2K’s electron-(anti)neutrino appearance, muon-(anti)neutrino disappearance, and joint three-flavour analyses, contributing to the majority of published T2K oscillation results.

In 2019, a major code refactoring effort culminated in the development of a generic, CPU-efficient, and highly flexible oscillation analysis framework: the VALOR Software Development Kit (SDK). This SDK was carved from the original T2K analysis, transforming it into a reusable, experiment-agnostic toolkit that eliminates the need for maintaining multiple forks.

Today, the VALOR SDK underpins advanced oscillation and sensitivity analyses for a range of current and next-generation neutrino experiments. These include T2K, JUNO, SBN, DUNE, and Hyper-Kamiokande, where each implementation adds only a thin experiment-specific interface layer atop the shared core framework.

This document serves as the official manual for the VALOR SDK.
